\documentclass[11pt, a4paper]{article}

% --- UNIVERSAL PREAMBLE BLOCK ---
\usepackage[a4paper, top=2.5cm, bottom=2.5cm, left=2cm, right=2cm]{geometry}
\usepackage{fontspec}

\usepackage[english, bidi=basic, provide=*]{babel}

\babelprovide[import, onchar=ids fonts]{english}

% Set default/Latin font to Sans Serif in the main (rm) slot
\babelfont{rm}{Noto Sans}

% --- ADDITIONAL PACKAGES ---
\usepackage{amssymb}    % For checkboxes \square
\usepackage{booktabs}   % For professional tables
\usepackage{longtable}  % For multi-page tracking table
\usepackage{enumitem}   % For customized lists
\usepackage{xcolor}     % For colored highlights
\usepackage[colorlinks=true, linkcolor=blue]{hyperref} % For internal links

% --- CUSTOM COMMANDS ---
\newcommand{\project}[1]{\vspace{1mm}\noindent\textbf{\textcolor{blue}{Hands-on Project:}} #1}
\newcommand{\usecase}[1]{\vspace{1mm}\noindent\textbf{\textcolor{teal}{Use Case \& Architecture:}} #1}
\newcommand{\capstone}[2]{
    \vspace{3mm}
    \noindent\fbox{
        \parbox{0.97\textwidth}{
            \textbf{\textcolor{red}{CAPSTONE PROJECT - #1:}} #2
        }
    }
    \vspace{3mm}
}
\newcommand{\dayitem}[2]{\item \textbf{#1:} #2}

\title{\Huge \textbf{Microsoft Fabric Master Data Engineer}\\ \Large 6-Month Exhaustive Training \& Certification Program}
\author{Advanced Architecture, Hands-on Labs, and MLOps}
\date{}

\begin{document}

\maketitle

\begin{abstract}
\noindent This 24-week program is designed to transform you from a practitioner to a Master Data Engineer on Microsoft Fabric. It progressively builds from OneLake foundations to advanced Medallion architectures, integrating daily theory, weekly service-specific mini-projects, architectural use case analyses, and massive monthly Capstone projects. A tracking timetable is included at the end.
\end{abstract}

\tableofcontents
\newpage

% ==========================================
% MILESTONE 1
% ==========================================
\section{Milestone 1: Core Data Infrastructure \& OneLake (Month 1)}
\textit{Objective: Master OneLake, Workspaces, Lakehouses, and the Medallion Architecture.}

\subsection*{Week 1: OneLake \& Fabric Foundations}
\begin{itemize}[label=-]
    \dayitem{Monday}{Introduction to Microsoft Fabric: SaaS architecture, Capacities, and Domains.}
    \dayitem{Tuesday}{OneLake concepts: The "OneDrive for Data", Workspaces, and Folders.}
    \dayitem{Wednesday}{OneLake Shortcuts: Virtualizing data from ADLS Gen2, Amazon S3, and Dataverse.}
    \dayitem{Thursday}{\project{Create a Fabric Workspace. Provision a Lakehouse and create a OneLake shortcut to a public Azure Blob Storage container without moving any data.}}
    \dayitem{Friday}{\usecase{Architect a multi-cloud data virtualization strategy for a global retailer that has historical sales in AWS S3 and current CRM data in Dataverse.}}
\end{itemize}

\subsection*{Week 2: The Medallion Architecture}
\begin{itemize}[label=-]
    \dayitem{Monday}{Medallion principles: Bronze (Raw), Silver (Validated), and Gold (Enriched).}
    \dayitem{Tuesday}{Lakehouse vs. Data Warehouse in Fabric: When to choose which.}
    \dayitem{Wednesday}{Organizing workspaces and Lakehouses for Medallion scaling.}
    \dayitem{Thursday}{\project{Deploy three Lakehouses (Bronze, Silver, Gold). Manually upload raw CSV files to Bronze, and set up the folder structure for future transformations.}}
    \dayitem{Friday}{\usecase{Design a multi-workspace architecture for a financial institution ensuring strict isolation between raw ingested logs and curated financial reports.}}
\end{itemize}

\subsection*{Week 3: Apache Spark in Fabric}
\begin{itemize}[label=-]
    \dayitem{Monday}{Spark compute in Fabric: Starter pools vs. Custom pools, Spark environments.}
    \dayitem{Tuesday}{Fabric Notebooks: PySpark basics, reading/writing files in OneLake.}
    \dayitem{Wednesday}{Spark SQL and combining PySpark with SQL in the same notebook.}
    \dayitem{Thursday}{\project{Write a PySpark notebook that reads raw CSV data from the Bronze Lakehouse, drops duplicates, drops nulls, and writes to the Silver Lakehouse as Delta tables.}}
    \dayitem{Friday}{\usecase{Design a scalable ingestion layer for a telemetry system processing 100 GB of JSON logs daily using Spark Notebooks.}}
\end{itemize}

\subsection*{Week 4: Delta Lake Optimization}
\begin{itemize}[label=-]
    \dayitem{Monday}{Delta Lake internals: Parquet files + Transaction logs (JSON).}
    \dayitem{Tuesday}{Optimizing Delta Tables: V-Order, OPTIMIZE, and Z-Ordering.}
    \dayitem{Wednesday}{Managing Delta Tables: VACUUM, Time Travel, and Schema Evolution.}
    \dayitem{Thursday}{\project{Run an OPTIMIZE command with V-Order on a large Silver table. Use Delta Time Travel to query a previous version of the table after an accidental deletion.}}
    \dayitem{Friday}{\usecase{Troubleshoot and re-architect a severely degraded Delta table that has accumulated millions of small files (small file problem) over 6 months of streaming.}}
\end{itemize}

\capstone{Milestone 1}{Build a Spark-driven Medallion Lakehouse. Create Workspaces for Bronze, Silver, and Gold. Use OneLake shortcuts to ingest sample data. Write parameterized Spark notebooks to clean the data (Silver) and aggregate it into a Star Schema (Gold) using Delta format with V-Order enabled.}

% ==========================================
% MILESTONE 2
% ==========================================
\section{Milestone 2: Data Warehousing \& BI Integration (Month 2)}
\textit{Objective: Deep dive into Synapse Data Warehouse, T-SQL, and DirectLake mode.}

\subsection*{Week 5: Synapse Data Warehouse in Fabric}
\begin{itemize}[label=-]
    \dayitem{Monday}{Fabric Data Warehouse architecture: Compute and Storage separation.}
    \dayitem{Tuesday}{T-SQL in Fabric: Supported surface area, COPY INTO, and cross-database querying.}
    \dayitem{Wednesday}{Data Warehouse ingestion patterns (Pipelines vs. T-SQL vs. Dataflows).}
    \dayitem{Thursday}{\project{Provision a Data Warehouse. Use the COPY INTO statement via a SQL Query to load 5 million rows of Parquet data directly from a OneLake path.}}
    \dayitem{Friday}{\usecase{Estimate capacity usage and recommend an ingestion pattern (Lakehouse vs. Warehouse) for a legacy SQL Server migration requiring heavy stored procedure usage.}}
\end{itemize}

\subsection*{Week 6: SQL Analytics Endpoint}
\begin{itemize}[label=-]
    \dayitem{Monday}{The automatic SQL Analytics Endpoint for Lakehouses (Read-only).}
    \dayitem{Tuesday}{Creating Views, TVFs, and Stored Procedures over Lakehouse Delta tables.}
    \dayitem{Wednesday}{Security: Object-level and Row-level security (RLS) in the SQL Endpoint.}
    \dayitem{Thursday}{\project{Connect to a Lakehouse's SQL Analytics Endpoint via SQL Server Management Studio (SSMS). Create a View joining a Silver Delta table with a Gold Delta table.}}
    \dayitem{Friday}{\usecase{Design a logical data warehouse architecture allowing data analysts to query raw data lake files using purely T-SQL without duplicating the data.}}
\end{itemize}

\subsection*{Week 7: Data Modeling for Power BI}
\begin{itemize}[label=-]
    \dayitem{Monday}{Dimensional modeling: Star schemas, Facts, and Dimensions.}
    \dayitem{Tuesday}{The default Power BI semantic model in Fabric Workspaces.}
    \dayitem{Wednesday}{Creating custom Semantic Models and defining relationships.}
    \dayitem{Thursday}{\project{Open the Semantic Model view in a Fabric Workspace. Define a 1-to-many relationship between a Date dimension and a Sales fact table. Create a DAX measure.}}
    \dayitem{Friday}{\usecase{Resolve a circular dependency issue in a complex enterprise data model by implementing a proper bridge table architecture.}}
\end{itemize}

\subsection*{Week 8: Power BI DirectLake Mode}
\begin{itemize}[label=-]
    \dayitem{Monday}{DirectLake concept: Reading Delta/Parquet directly into memory (No Import, No DirectQuery).}
    \dayitem{Tuesday}{DirectLake fallback behavior (falling back to DirectQuery).}
    \dayitem{Wednesday}{Optimizing models for DirectLake (V-Order importance).}
    \dayitem{Thursday}{\project{Create a Power BI report using DirectLake mode against your Gold Lakehouse. Benchmark the performance vs a traditional DirectQuery model.}}
    \dayitem{Friday}{\usecase{Architect a real-time BI dashboard for an executive team that requires sub-second response times on 100-million-row datasets without nightly refresh delays.}}
\end{itemize}

\capstone{Milestone 2}{End-to-End Analytics Workflow. Establish a Synapse Data Warehouse. Create complex T-SQL stored procedures to build a Star Schema from raw OneLake data. Build a custom Semantic Model with relationships and DAX measures, and publish a DirectLake Power BI report.}

% ==========================================
% MILESTONE 3
% ==========================================
\section{Milestone 3: Batch Processing \& ETL Orchestration (Month 3)}
\textit{Objective: Master Data Factory in Fabric, Dataflows Gen2, and Pipelines.}

\subsection*{Week 9: Data Factory Pipelines}
\begin{itemize}[label=-]
    \dayitem{Monday}{Fabric Pipelines: Copy Activity, Lookups, and Get Metadata.}
    \dayitem{Tuesday}{Control Flow: ForEach, If Condition, Until, and Switch.}
    \dayitem{Wednesday}{Parameters, Variables, and dynamic expressions in Pipelines.}
    \dayitem{Thursday}{\project{Build a dynamic pipeline that looks up a list of table names from a configuration file and uses a ForEach loop to copy them from an Azure SQL DB to OneLake.}}
    \dayitem{Friday}{\usecase{Design a metadata-driven ingestion framework capable of onboarding hundreds of new tables automatically without manual pipeline creation.}}
\end{itemize}

\subsection*{Week 10: Dataflows Gen2}
\begin{itemize}[label=-]
    \dayitem{Monday}{Power Query Online integration: Visual data transformations.}
    \dayitem{Tuesday}{Dataflows Gen2 output destinations (Lakehouse, Warehouse, Azure SQL).}
    \dayitem{Wednesday}{Staging vs. Data destinations, incremental refresh in Dataflows.}
    \dayitem{Thursday}{\project{Build a Dataflow Gen2 to connect to a web API, pivot and clean the JSON response using Power Query, and set the destination to a Gold Lakehouse table.}}
    \dayitem{Friday}{\usecase{Evaluate Mapping Dataflows (ADF), Dataflows Gen2, and Spark Notebooks for a team of business analysts migrating from Excel-based ETL.}}
\end{itemize}

\subsection*{Week 11: Data Integration \& Connectivity}
\begin{itemize}[label=-]
    \dayitem{Monday}{On-premises Data Gateway vs. VNet Data Gateway.}
    \dayitem{Tuesday}{Connecting to hybrid sources (SQL Server, Oracle, SAP).}
    \dayitem{Wednesday}{Integrating REST APIs and Webhooks.}
    \dayitem{Thursday}{\project{Configure a mock on-premises SQL Server. Install the On-premises Data Gateway and build a Fabric pipeline to extract data daily to OneLake.}}
    \dayitem{Friday}{\usecase{Architect a secure hybrid connectivity pattern for a bank that strictly forbids public IP endpoints for their on-premise transactional databases.}}
\end{itemize}

\subsection*{Week 12: Orchestration \& Monitoring}
\begin{itemize}[label=-]
    \dayitem{Monday}{Invoking Notebooks and Dataflows from Pipelines.}
    \dayitem{Tuesday}{Dependency management (On Success, On Failure, On Completion).}
    \dayitem{Wednesday}{Monitoring Hub, alert routing, and Pipeline run history.}
    \dayitem{Thursday}{\project{Create a master orchestration pipeline that triggers a Dataflow Gen2, waits for completion, triggers a Spark Notebook, and sends a Teams/Outlook alert on failure.}}
    \dayitem{Friday}{\usecase{Design a resilient, self-healing orchestration layer for a daily banking reconciliation process spanning 15 dependent systems using Fabric Pipelines.}}
\end{itemize}

\capstone{Milestone 3}{Automated Batch ELT Pipeline. Use Fabric Data Factory to build a complete pipeline. Ingest data via REST APIs using the Copy Activity, use Dataflows Gen2 for visual data cleansing, and orchestrate Spark notebooks for complex ML feature engineering. Schedule the pipeline and configure failure alerts.}

\newpage
% ==========================================
% MILESTONE 4
% ==========================================
\section{Milestone 4: Real-Time Intelligence (Month 4)}
\textit{Objective: Conquer unbounded data, Eventstreams, and KQL Databases.}

\subsection*{Week 13: Eventstreams \& Real-Time Ingestion}
\begin{itemize}[label=-]
    \dayitem{Monday}{Real-Time Intelligence architecture in Fabric.}
    \dayitem{Tuesday}{Eventstreams: Sources (Azure Event Hubs, Custom Apps, IoT).}
    \dayitem{Wednesday}{Eventstream Destinations (Lakehouse, KQL Database, Reflex).}
    \dayitem{Thursday}{\project{Create a Fabric Eventstream. Write a Python script to send mock IoT temperature data to the custom app endpoint, and route the data to a KQL Database.}}
    \dayitem{Friday}{\usecase{Architect a global telemetry ingestion point for 5 million mobile games, ensuring high throughput and low-latency processing into Fabric.}}
\end{itemize}

\subsection*{Week 14: KQL Databases (Eventhouses)}
\begin{itemize}[label=-]
    \dayitem{Monday}{Eventhouses and Kusto Query Language (KQL) databases.}
    \dayitem{Tuesday}{KQL syntax basics: search, where, project, summarize.}
    \dayitem{Wednesday}{Advanced KQL: Time series analysis, joins, and windowing.}
    \dayitem{Thursday}{\project{Query the IoT data in your KQL database. Write a KQL query to calculate the 5-minute rolling average temperature and detect anomalies using built-in functions.}}
    \dayitem{Friday}{\usecase{Troubleshoot a scenario where a streaming dashboard is showing inaccurate daily aggregations due to mobile devices going offline and uploading data late.}}
\end{itemize}

\subsection*{Week 15: Data Activator (Reflex)}
\begin{itemize}[label=-]
    \dayitem{Monday}{Data Activator concepts: Triggers, Conditions, and Actions.}
    \dayitem{Tuesday}{Creating Reflex items from Power BI visuals or Eventstreams.}
    \dayitem{Wednesday}{Automated actions: Teams messages, Emails, and custom workflows.}
    \dayitem{Thursday}{\project{Set up a Data Activator trigger on your streaming IoT data. Configure an alert to send an email if the temperature exceeds a threshold for 3 consecutive minutes.}}
    \dayitem{Friday}{\usecase{Design a proactive alerting system for a logistics company to instantly notify drivers via custom app webhooks when a delivery truck goes off route.}}
\end{itemize}

\subsection*{Week 16: Real-Time Dashboards}
\begin{itemize}[label=-]
    \dayitem{Monday}{KQL Querysets vs. Real-Time Dashboards.}
    \dayitem{Tuesday}{Building visual dashboards directly on KQL databases.}
    \dayitem{Wednesday}{Power BI automatic page refresh (DirectQuery vs KQL).}
    \dayitem{Thursday}{\project{Create a Real-Time Dashboard natively in Fabric using KQL queries. Add auto-refreshing charts for active IoT devices and current temperature trends.}}
    \dayitem{Friday}{\usecase{Design the BI access layer for a manufacturing floor where operators need sub-second latency dashboards monitoring machine health (vibration and heat).}}
\end{itemize}

\capstone{Milestone 4}{Real-Time Fraud Detection. Build a Python script simulating credit card swipes sent to a Fabric Eventstream. Route data to a KQL Database. Write KQL queries to detect if a card is used in two different cities within 10 minutes. Use Data Activator to instantly alert the security team on Microsoft Teams.}


% ==========================================
% MILESTONE 5
% ==========================================
\section{Milestone 5: Data Science \& AI Operationalization (Month 5)}
\textit{Objective: Embed AI into data pipelines using MLflow and SynapseML.}

\subsection*{Week 17: Data Science Foundations in Fabric}
\begin{itemize}[label=-]
    \dayitem{Monday}{The Data Science experience: Notebooks, Data Wrangler.}
    \dayitem{Tuesday}{Exploratory Data Analysis (EDA) with Pandas and PySpark.}
    \dayitem{Wednesday}{Feature engineering and preparation for machine learning.}
    \dayitem{Thursday}{\project{Use Fabric Data Wrangler in a Notebook to visually clean a dataset, handle missing values, and generate the corresponding PySpark code automatically.}}
    \dayitem{Friday}{\usecase{Architect a system to automatically forecast daily inventory needs for 1,000 retail stores using purely Spark-based pipelines.}}
\end{itemize}

\subsection*{Week 18: Model Training \& MLflow}
\begin{itemize}[label=-]
    \dayitem{Monday}{Machine Learning Experiments and Models in Fabric.}
    \dayitem{Tuesday}{Using MLflow for tracking parameters, metrics, and artifacts.}
    \dayitem{Wednesday}{Registering models in the Fabric Model Registry.}
    \dayitem{Thursday}{\project{Train a Scikit-Learn Random Forest model. Use `mlflow.autolog()` to automatically track the experiment in Fabric, and register the best model.}}
    \dayitem{Friday}{\usecase{Design a workflow for data scientists to track hundreds of hyperparameter tuning runs collaboratively within a shared Fabric Workspace.}}
\end{itemize}

\subsection*{Week 19: SynapseML \& AI Integration}
\begin{itemize}[label=-]
    \dayitem{Monday}{Introduction to SynapseML (Distributed ML library).}
    \dayitem{Tuesday}{Pre-built AI Models: Integrating Azure OpenAI and Cognitive Services.}
    \dayitem{Wednesday}{Copilot in Fabric: Generating code and insights using natural language.}
    \dayitem{Thursday}{\project{Use SynapseML in a Spark notebook to call the Azure OpenAI API. Pass a column of customer reviews and append a new column with the sentiment analysis result.}}
    \dayitem{Friday}{\usecase{Design a scalable pipeline to process terabytes of uploaded user reviews, extract sentiment, and make the metadata available in a Gold Lakehouse.}}
\end{itemize}

\subsection*{Week 20: MLOps \& Batch Scoring}
\begin{itemize}[label=-]
    \dayitem{Monday}{Batch scoring with Spark `PREDICT` function.}
    \dayitem{Tuesday}{Applying registered MLflow models to Delta tables.}
    \dayitem{Wednesday}{Automating model inference via Fabric Pipelines.}
    \dayitem{Thursday}{\project{Load a registered MLflow model and use the `PREDICT` function in Spark to score a new batch of 100,000 rows. Write the predictions back to OneLake.}}
    \dayitem{Friday}{\usecase{Architect an MLOps lifecycle that detects data drift in a production recommendation engine and automatically triggers a retraining Pipeline.}}
\end{itemize}

\capstone{Milestone 5}{End-to-End MLOps Pipeline. Use Data Factory Pipelines to orchestrate a workflow that extracts features from a Lakehouse, triggers a Notebook to train a churn prediction model, tracks it with MLflow, registers the model, and runs a daily batch scoring Notebook that populates a Power BI DirectLake report.}

\newpage
% ==========================================
% MILESTONE 6
% ==========================================
\section{Milestone 6: Security, Administration \& Exam Prep (Month 6)}
\textit{Objective: Secure OneLake, implement Purview, and prepare for DP-600.}

\subsection*{Week 21: Security \& Access Control}
\begin{itemize}[label=-]
    \dayitem{Monday}{Workspace Roles: Admin, Member, Contributor, Viewer.}
    \dayitem{Tuesday}{Item-level permissions vs. Workspace-level permissions.}
    \dayitem{Wednesday}{OneLake Data Access Roles (Preview) and Row-Level Security (RLS) in SQL.}
    \dayitem{Thursday}{\project{Create a Workspace and assign a Contributor role. Configure Row-Level Security on a SQL Analytics Endpoint to restrict a user to seeing only 'US' region data.}}
    \dayitem{Friday}{\usecase{Design a highly secure data perimeter for a healthcare provider ensuring strict separation of duties between Data Engineers and Data Analysts within Fabric.}}
\end{itemize}

\subsection*{Week 22: Governance \& Microsoft Purview in Fabric}
\begin{itemize}[label=-]
    \dayitem{Monday}{Purview Hub in Fabric and tenant-level administration.}
    \dayitem{Tuesday}{Information Protection: Sensitivity labels applied to Fabric items.}
    \dayitem{Wednesday}{Data Lineage, Endorsement (Promoted/Certified), and Impact Analysis.}
    \dayitem{Thursday}{\project{Apply a 'Highly Confidential' sensitivity label to a Lakehouse. View the automated data lineage graph from the Lakehouse through a Pipeline to a Power BI report.}}
    \dayitem{Friday}{\usecase{Architect a Data Mesh for a global conglomerate with 5 independent domains, ensuring centralized governance via Purview while using separate Fabric Workspaces.}}
\end{itemize}

\subsection*{Week 23: Administration \& CI/CD}
\begin{itemize}[label=-]
    \dayitem{Monday}{Fabric Capacity management, metrics app, and pausing/resuming.}
    \dayitem{Tuesday}{Git Integration: Syncing Workspaces with Azure DevOps/GitHub.}
    \dayitem{Wednesday}{Deployment Pipelines: Dev $\rightarrow$ Test $\rightarrow$ Prod promotion.}
    \dayitem{Thursday}{\project{Connect your Workspace to an Azure DevOps repository. Commit your Notebooks and Pipelines. Create a Deployment Pipeline to promote items to a 'Prod' workspace.}}
    \dayitem{Friday}{\usecase{Design an enterprise CI/CD pipeline allowing multiple data engineering teams to collaborate on the same Lakehouse without overwriting each other's code.}}
\end{itemize}

\subsection*{Week 24: DP-600 Certification Preparation}
\begin{itemize}[label=-]
    \dayitem{Monday}{Review DP-600 (Implementing Analytics Solutions Using Microsoft Fabric) Objectives.}
    \dayitem{Tuesday}{Deep dive: DirectLake vs Import vs DirectQuery optimization.}
    \dayitem{Wednesday}{Full-length Mock Exam 1 (Identify weak domains).}
    \dayitem{Thursday}{Targeted review of weak domains (e.g., Spark optimization, Pipeline parameters).}
    \dayitem{Friday}{Full-length Mock Exam 2.}
\end{itemize}

\capstone{Milestone 6}{Secure, Governed Data Platform. Build a multi-workspace architecture. Ingest streaming PII data. Use Dataflows Gen2 to mask the PII. Apply Sensitivity Labels and Endorsements. Connect to Azure DevOps and use Deployment Pipelines to seamlessly promote the architecture from a Development Workspace to Production.}

\newpage
% ==========================================
% TRACKING TIMETABLE
% ==========================================
\section{6-Month Progress Tracking Timetable}

\begin{center}
\renewcommand{\arraystretch}{1.5}
\begin{longtable}{@{} c l p{7cm} c @{}}
\toprule
\textbf{Week} & \textbf{Primary Focus} & \textbf{Key Project / Capstone} & \textbf{Status} \\ 
\midrule
\endfirsthead

\toprule
\textbf{Week} & \textbf{Primary Focus} & \textbf{Key Project / Capstone} & \textbf{Status} \\ 
\midrule
\endhead

\bottomrule
\endfoot

\bottomrule
\endlastfoot

1 & OneLake \& Fabric Foundations & OneLake Shortcuts & $\square$ \\
2 & The Medallion Architecture & Bronze/Silver/Gold Lakehouses & $\square$ \\
3 & Apache Spark in Fabric & PySpark Data Cleansing & $\square$ \\
4 & Delta Lake Optimization & \textbf{Cap 1: Spark Medallion Lakehouse} & $\square$ \\
\midrule
5 & Synapse Data Warehouse & COPY INTO from OneLake & $\square$ \\
6 & SQL Analytics Endpoint & Cross-Lakehouse T-SQL Views & $\square$ \\
7 & Data Modeling for Power BI & Fabric Semantic Models \& DAX & $\square$ \\
8 & Power BI DirectLake Mode & \textbf{Cap 2: End-to-End Analytics Workflow} & $\square$ \\
\midrule
9 & Data Factory Pipelines & Dynamic Pipeline with ForEach Loop & $\square$ \\
10 & Dataflows Gen2 & Power Query Online to Lakehouse & $\square$ \\
11 & Connectivity \& Hybrid & On-Premises Data Gateway config & $\square$ \\
12 & Orchestration & \textbf{Cap 3: Automated Batch ELT} & $\square$ \\
\midrule
13 & Eventstreams & Custom App Eventstream Ingestion & $\square$ \\
14 & KQL Databases & KQL Queries \& Anomaly Detection & $\square$ \\
15 & Data Activator (Reflex) & Automated Alerts via Teams/Email & $\square$ \\
16 & Real-Time Dashboards & \textbf{Cap 4: Real-Time Fraud Telemetry} & $\square$ \\
\midrule
17 & Data Science Foundations & Data Wrangler code generation & $\square$ \\
18 & Model Training \& MLflow & MLflow Autologging \& Registration & $\square$ \\
19 & SynapseML \& AI Integration & Azure OpenAI sentiment analysis & $\square$ \\
20 & MLOps \& Batch Scoring & \textbf{Cap 5: End-to-End MLOps Pipeline} & $\square$ \\
\midrule
21 & Security \& Access Control & SQL RLS and Workspace Roles & $\square$ \\
22 & Governance \& Purview & Sensitivity Labels \& Lineage Hub & $\square$ \\
23 & Administration \& CI/CD & Git Integration \& Deployment Pipelines & $\square$ \\
24 & Exam Prep (DP-600) & \textbf{Cap 6: Secure Governed Platform} & $\square$ \\

\end{longtable}
\end{center}

\vspace{1cm}
\begin{center}
    \textit{"Consistency is the architecture of mastery."} \\
    Best of luck on your 6-month journey to becoming a Master Microsoft Fabric Analytics Engineer!
\end{center}

\end{document}